% =========================================================
% SECCIÓN 3.4 - FORMULACIÓN DEL SISTEMA DE ECUACIONES DIFERENCIALES
% =========================================================


Con las variables de estado y el dominio físico establecidos en la Sección \ref{sec:cap3_variables}, esta sección formula el sistema de ecuaciones diferenciales ordinarias (EDO) que gobierna la evolución del vector $\mathbf{x}(t)$ en el reactor batch. El sistema se deriva directamente de los flujos metabólicos definidos en la Sección \ref{subsec:cap3_flujos}, imponiendo la conservación de masa de carbono en cada compartimento. Las formas explícitas de las funciones cinéticas —ingesta, asimilación y modulación térmica— se detallan en la Sección \ref{sec:cap3_cinetica}.

Se asume que el sistema satisface las condiciones de Carathéodory locales en $\Omega$, garantizando existencia y unicidad de la solución para condiciones iniciales $\mathbf{x}_0 \in \Omega$; las demostraciones formales se remiten al Anexo \ref{anexo:formulacion_matematica} para no interrumpir el desarrollo del capítulo.

% ---------------------------------------------------------
\subsection{Balance del sustrato orgánico}
\label{subsec:cap3_ode_sustrato}
% ---------------------------------------------------------

La masa de carbono orgánico disponible $S(t)$ disminuye exclusivamente por la ingesta larval, dado que la degradación microbiana del sustrato no ingerido se modela como una fuente de \ch{CO2} externa al balance de sustrato (vía $\dot{D}_{\mu}$ en la Sección \ref{subsec:cap3_ode_co2}). La ecuación de balance es:
\begin{equation}
    \frac{\mathrm{d}S}{\mathrm{d}t} = -\dot{I}_{\mathrm{tot}}(t) = -N_{\mathrm{larvas}} \cdot \tilde{\imath}\bigl(S(t), B(t), T(t)\bigr),
    \label{eq:cap3_ode_sustrato}
\end{equation}
donde $\tilde{\imath}$ es la tasa de ingesta específica por larva (g C/larva/min), función creciente de la biomasa estructural $B$ —que determina la superficie bucal y la capacidad de procesamiento— y de la disponibilidad de sustrato $S$. La dependencia térmica se incorpora mediante el factor $f(T)$ definido en la Sección \ref{sec:cap3_cinetica}. La forma funcional de $\tilde{\imath}$ sigue la cinética de Holling II, cuya parametrización se presenta en la Sección \ref{subsec:cap3_holling}.

La condición inicial $S(0) = S_0$ está dada por la ecuación \eqref{eq:cap3_S0}. El dominio requiere $S(t) \geq 0$; cuando $S \to 0$, la ingesta se anula ($\tilde{\imath} \to 0$) y el sustrato permanece no negativo, comportamiento que se verifica cualitativamente en la Sección \ref{sec:cap3_analisis}.

% ---------------------------------------------------------
\subsection{Balance de asimilados en estado cuasi-estacionario}
\label{subsec:cap3_ode_reserva}
% ---------------------------------------------------------

La reserva de asimilados $E(t)$ opera como un compartimento de amortiguamiento de alta velocidad respecto a la dinámica de crecimiento. Siguiendo la hipótesis de separación de escalas temporales propia del marco DEB \parencite{eriksenDynamicModellingFeed2022}, se impone el estado cuasi-estacionario:
\begin{equation}
    \frac{\mathrm{d}E}{\mathrm{d}t} \approx 0,
    \label{eq:cap3_qss_reserva}
\end{equation}
lo cual equivale a igualar el flujo de entrada (asimilación) con los flujos de salida (mantenimiento y síntesis). La asimilación total se expresa como $\dot{A} = \kappa_A \cdot \dot{I}_{\mathrm{tot}}$, donde $\kappa_A \in (0,1]$ es la eficiencia de asimilación intestinal. El balance resultante es:
\begin{equation}
    \dot{A}(t) = \dot{M}(t) + \frac{\dot{G}_B(t)}{Y_B} + \frac{\dot{G}_L(t)}{Y_L},
    \label{eq:cap3_balance_asimilados}
\end{equation}
donde:
\begin{align}
    \dot{M}(t) &= m \cdot B(t) \cdot f\bigl(T(t)\bigr)
    \label{eq:cap3_mantenimiento}
\end{align}
es la tasa de mantenimiento somático (g C/min), con $m$ la tasa específica de mantenimiento a la temperatura de referencia; $\dot{G}_B(t)$ y $\dot{G}_L(t)$ son los flujos de síntesis de biomasa estructural y lípidos (g C/min); $Y_B, Y_L \in (0,1]$ son los rendimientos de conversión de reserva en biomasa y lípidos, respectivamente, de modo que $\dot{G}_B/Y_B$ representa el carbono total extraído de la reserva incluyendo la fracción disipada como \ch{CO2} durante la síntesis (\emph{overhead} metabólico).

La partición del excedente asimilado neto $\dot{A} - \dot{M}$ entre los compartimentos de crecimiento obedece a la regla de prioridad $\kappa$ del DEB. Si $\dot{A} > \dot{M}$:
\begin{equation}
    \dot{G}_B(t) + \dot{G}_L(t) = \kappa \cdot \bigl[\dot{A}(t) - \dot{M}(t)\bigr],
    \label{eq:cap3_prioridad_crecimiento}
\end{equation}
y si $\dot{A} \leq \dot{M}$:
\begin{equation}
    \dot{G}_B(t) = \dot{G}_L(t) = 0,
    \label{eq:cap3_prioridad_mantenimiento}
\end{equation}
donde $\kappa \in (0,1]$ es el parámetro de prioridad de crecimiento versus mantenimiento. La asignación específica entre $\dot{G}_B$ y $\dot{G}_L$ —esto es, la fracción del flujo de crecimiento destinada a lípidos en función del estadio larval— sigue la parametrización etaria de \parencite{eriksenDynamicModellingFeed2022}, documentada en el Anexo \ref{anexo:formulacion_matematica}.

% ---------------------------------------------------------
\subsection{Dinámica de la biomasa estructural}
\label{subsec:cap3_ode_biomasa}
% ---------------------------------------------------------

La biomasa estructural $B(t)$ crece por deposición de carbono desde la reserva y puede decrecer por removilización catabólica en condiciones de ayuno extremo:
\begin{equation}
    \frac{\mathrm{d}B}{\mathrm{d}t} = \dot{G}_B(t) - \delta_B \cdot B(t),
    \label{eq:cap3_ode_biomasa_general}
\end{equation}
donde $\delta_B$ es la tasa específica de removilización estructural. En el régimen experimental del Capítulo \ref{chap:instrumentacion} —alimentación \emph{ad libitum} con renovación cada 48 h— la condición $\dot{A} > \dot{M}$ se satisface durante la mayor parte del ciclo, por lo que $\dot{G}_B > 0$ y $\delta_B \cdot B \approx 0$. La ecuación se reduce entonces a:
\begin{equation}
    \frac{\mathrm{d}B}{\mathrm{d}t} \approx \dot{G}_B(t),
    \label{eq:cap3_ode_biomasa}
\end{equation}
que vincula directamente la curva de crecimiento teórica con las mediciones destructivas de biomasa húmeda. La condición inicial $B(0) = B_0$ proviene de la ecuación \eqref{eq:cap3_B0}.

% ---------------------------------------------------------
\subsection{Acumulación de lípidos de reserva}
\label{subsec:cap3_ode_lipidos}
% ---------------------------------------------------------

El compartimento lipídico $L(t)$ acumula carbono destinado al almacenamiento energético de largo plazo, predominante en estadios tardíos previos a la metamorfosis:
\begin{equation}
    \frac{\mathrm{d}L}{\mathrm{d}t} = \dot{G}_L(t) - \delta_L \cdot L(t),
    \label{eq:cap3_ode_lipidos_general}
\end{equation}
con $\delta_L$ la tasa de movilización lipídica. Bajo el mismo argumento de alimentación suficiente, $\delta_L \approx 0$ y la dinámica se simplifica a $\mathrm{d}L/\mathrm{d}t \approx \dot{G}_L(t)$. La condición inicial $L(0) = L_0 = 0$ refleja la ausencia de reservas significativas en estadio L3 inicial.

La acumulación de lípidos constituye un sumidero de carbono paralelo a la biomasa estructural. Su magnitud relativa —que puede alcanzar hasta el $35\%$ de la biomasa seca en sustratos ricos en carbohidratos según \textcite{grausaMetabolicModelingHermetia2023}— condiciona directamente el balance de carbono disponible para emisión como \ch{CO2}, motivando su inclusión explícita en el modelo en lugar de agruparla en un único compartimento de ``biomasa total''.

% ---------------------------------------------------------
\subsection{Producción y acumulación de \texorpdfstring{\ch{CO2}}{CO2}}
\label{subsec:cap3_ode_co2}
% ---------------------------------------------------------

La quinta ecuación del sistema acumula el carbono emitido como dióxido de carbono, integrando las fuentes larval y microbiana:
\begin{equation}
    \frac{\mathrm{d}C_{\ch{CO2}}}{\mathrm{d}t} = \underbrace{\dot{M}(t) + \bigl(1-Y_B\bigr)\dot{G}_B(t) + \bigl(1-Y_L\bigr)\dot{G}_L(t)}_{\dot{D}_{\mathrm{larval}}(t)} + \underbrace{\dot{D}_{\mu}(t)}_{\mathrm{microbiana}},
    \label{eq:cap3_ode_co2}
\end{equation}
donde $\dot{D}_{\mathrm{larval}}$ es la disipación metabólica directa de las larvas. Los términos $(1-Y_B)\dot{G}_B$ y $(1-Y_L)\dot{G}_L$ capturan el costo de síntesis (\emph{overhead}): la fracción de reserva consumida que no se incorpora a tejido sino que se pierde como \ch{CO2} en las reacciones anabólicas. El término $\dot{D}_{\mu}(t)$ modela la respiración aeróbica de la microbiota asociada al sustrato no ingerido o al tracto digestivo, que se trata como una fuente externa no acoplada al metabolismo larval en las presentes EDO.

La conexión con la señal del sensor NDIR se establece mediante la conversión de moles de carbono a gramos de \ch{CO2} por minuto:
\begin{equation}
    \dot{m}_{\ch{CO2}}(t) = \frac{M_{\ch{CO2}}}{M_{\ch{C}}} \cdot \frac{\mathrm{d}C_{\ch{CO2}}}{\mathrm{d}t},
    \label{eq:cap3_conversion_masica}
\end{equation}
donde $M_{\ch{CO2}} = 44.01$ g/mol y $M_{\ch{C}} = 12.01$ g/mol. Esta expresión es consistente con la ecuación \eqref{eq:cap3_puente_sensor} y permite confrontar directamente la predicción del modelo con las pendientes $\Delta C/\Delta t$ convertidas a flujo másico en el protocolo de ventanas cerradas del Capítulo \ref{chap:instrumentacion}.

Es importante recalcar que el metano (\ch{CH4}) no aparece en las ecuaciones \eqref{eq:cap3_ode_co2} ni \eqref{eq:cap3_conversion_masica}. Dado que el modelo asume metabolismo aeróbico exclusivo, cualquier emisión de \ch{CH4} detectada por el sensor MQ-4 se interpreta como una desviación del régimen ideal, cuya detección y cuantificación se abordan en el Capítulo \ref{chap:metricas} como variable residual de diagnóstico operativo.

En conjunto, las ecuaciones \eqref{eq:cap3_ode_sustrato}--\eqref{eq:cap3_ode_co2} constituyen un sistema de cinco EDO acopladas no lineales, gobernado por los parámetros semilla $\{\kappa_A, m, a_{\max}, Y_B, Y_L, \kappa\}$ y las variables libres de calibración $\{\tau_h, \beta_0, \rho_{\mathrm{conv}}, T_A, T_{\mathrm{opt}}, T_H\}$. La siguiente sección detalla la cinética de asimilación y la modulación ambiental que cierran las formas funcionales de los flujos metabólicos.
